\section{Statement-adequacy ledger}\label{app:adequacy}

This appendix reproduces verbatim the formal statements on which the
manuscript places Lean-level trust.  It is intended to make statement
adequacy refereeable without requiring the reader to inspect the full
development.  The source anchors are
\path{RelativeConicArcs/PRSResidualQuadratic.lean} and
\path{RelativeConicArcs/PRSRedundancyNine.lean}; the import and audit
gates are those named in Section~\ref{sec:verification}.

\paragraph{Bottom definitions.}
\begingroup\footnotesize
\begin{verbatim}
def dividedPowerContraction {n : ℕ} (r : R)
    (a : Fin (n + 2) → R) : Fin (n + 1) → R :=
  fun i => a i.succ - r * a i.castSucc

structure ResidualSliceInput (S : Type*) (q : ℕ) where
  isDeep : S → Prop
  exceptional : S → Prop
  integralGenusAtMostOneSlice : S → Prop
  rationalPointOutsideDeletedDivisors :
    q ≥ 53 → S → Prop
  splitSquarefreeWitness : S → Prop
  witnessMakesShallow :
    ∀ {s}, splitSquarefreeWitness s → ¬ isDeep s
  pointGivesWitness :
    ∀ (hq : q ≥ 53) {s}, exceptional s →
      integralGenusAtMostOneSlice s →
      rationalPointOutsideDeletedDivisors hq s →
      splitSquarefreeWitness s

structure PersistentFamilyData (S : Type*) (q : ℕ)
    [Fintype S] [DecidableEq S] where
  tangent : Finset S
  sigma : Finset S
  deep : Finset S
  tangent_sigma_disjoint : Disjoint tangent sigma
  deep_eq : deep = tangent ∪ sigma
  fieldOrderPositive : 0 < q
  tangent_card : tangent.card = q * (q + 1)
  sigma_card_doubled : 2 * sigma.card = q * (q ^ 2 - 1)
  orbitCase : OrbitArithmeticCase
\end{verbatim}
\endgroup

The first definition is the paper's contraction.  The two structures
make geometric integrality, deleted-divisor rational points,
witness production, family identification, and family counts explicit
inputs rather than axioms hidden inside a Boolean certificate.

\paragraph{Headline synthesis terminal.}
\begingroup\footnotesize
\begin{verbatim}
theorem redundancyNineSynthesis {S : Type*} {q : ℕ}
    [Fintype S] [DecidableEq S]
    (hq : q ≥ 53)
    (slice : ResidualSliceInput S q)
    (families : PersistentFamilyData S q)
    (exceptional_iff :
      ∀ s, slice.exceptional s ↔ s ∉ families.deep)
    (persistentDeep :
      ∀ s, s ∈ families.deep → slice.isDeep s)
    (hslice :
      ∀ s, slice.exceptional s →
        slice.integralGenusAtMostOneSlice s)
    (hpoint :
      ∀ s (_hs : slice.exceptional s),
        slice.rationalPointOutsideDeletedDivisors hq s) :
    (∀ s, slice.isDeep s ↔ s ∈ families.deep) ∧
      (∀ s, s ∉ families.deep →
        slice.splitSquarefreeWitness s) ∧
      families.deep.card = q * (q + 1) ^ 2 / 2 ∧
      (projectiveOrbitCount families.orbitCase,
        semilinearOrbitCount families.orbitCase) =
      orbitCountPair families.orbitCase
\end{verbatim}
\endgroup

This terminal is adequate for the final synthesis implication once the paper supplies
the arguments \texttt{exceptional\_iff}, \texttt{persistentDeep},
\texttt{hslice}, and \texttt{hpoint}.  It does not prove those
geometric and coding inputs, nor does it derive the orbit table from a
formal group action.

\paragraph{Axiom audit.}
The guarded audit of fourteen terminals reports no axioms for the
finite witness and orbit-pair terminals; only
\texttt{propext} and \texttt{Quot.sound} for ring normalization; and
only \texttt{propext}, \texttt{Classical.choice}, and
\texttt{Quot.sound} for field/root terminals, the family cardinality,
and the headline synthesis theorem.  It uses no project-specific
axiom, native evaluator, finite-data import, or opaque computational
oracle.

